% Part: first-order-logic
% Chapter: models-theories
% Section: expressing-props-of-structures

\documentclass[../../../include/open-logic-section]{subfiles}

\begin{document}

\olfileid{fol}{mat}{exs}

\olsection{Nyatakake Sipat-Sipat \printtoken{P}{structure}}

\begin{explain}
Kerep migunani lan wigati kanggo nyatakake syarat tumrap
fungsi lan relasi, utawa luwih umum, yen fungsi lan relasi
ing sawijining struktur nyukupi syarat-syarat kasebut.
Upamane, kanggo sawijining basa kita kepengin
bisa mbedakake !!{structure}s kang nggambarake
teges kang kita karepake kanggo !!{predicate}s
saka struktur kang ora mangkono. Mesthi kita bebas nemtokake
!!{structure}s endi kang kita karepake, upamane,
kita bisa nemtokake yen interpretasi !!{predicate}~$\le$
kudu dadi relasi urutan, utawa yen kita mung nggatekake
interpretasi~$\Lang L$ kang ranahé dumadi saka himpunan
lan $\Obj \in$ diinterpretasi minangka relasi
"minangka !!a{element} saka". Nanging apa kita bisa
nindakake iki nganggo !!{sentence}s ing basa kasebut?
Kanthi tembung liya, syarat apa marang
!!a{structure}~$\Struct M$ kang bisa kita nyatakake
nganggo !!a{sentence} (utawa bisa uga himpunan
!!{sentence}s) ing basa~$\Struct M$? Ana sawetara
syarat kang ora bisa kita nyatakake. Upamane,
ora ana ukara ing~$\Lang L_A$ kang bisa dicukupi lan mung bener
ing sawijining !!{structure}~$\Struct M$ yen
$\Domain M = \Nat$. Cathetan koreksi sumber OLPL-751: watesan iki tumrap ukara kang bisa dicukupi, amarga ukara tanpa model marakake "mung yen" bener sacara vakum; salinan isomorfis bisa nduweni ranah liya. Kita ora bisa nyatakake
"ranahé mung ngemot wilangan asli". Nanging,
ana "sipat struktural" saka !!{structure}s
kang mbokmenawa bisa kita nyatakake. Sipat apa
saka !!{structure}s kang bisa kita nyatakake
nganggo !!{sentence}s? Utawa, kanthi cara liyane,
klumpukan !!{structure}s endi kang bisa kita
jlentrehake minangka struktur-struktur kang
ndadekake !!a{sentence} (utawa himpunan
!!{sentence}s) bener?
\end{explain}

\begin{defn}[Model saka himpunan]
Upamakna $\Gamma$ iku himpunan !!{sentence}s ing
basa~$\Lang L$. Kita kandha yen !!a{structure}~$\Struct M$
\emph{iku model saka}~$\Gamma$ yen $\Sat{M}{!A}$
kanggo saben $!A \in \Gamma$.
\end{defn}

\begin{ex}
Ukara $\lforall[x][x \le x]$ bener ing~$\Struct M$
yen lan mung yen $\Assign{\le}{M}$ iku relasi refleksif.
Ukara
$\lforall[x][\lforall[y][((x \le y \land y \le x) \lif x = y)]]$
bener ing~$\Struct M$ yen lan mung yen $\Assign{\le}{M}$
iku antisimetris. Ukara
$\lforall[x][\lforall[y][\lforall[z][((x \le y \land y \le z)
      \lif x \le z)]]]$ bener ing~$\Struct M$ yen lan
mung yen $\Assign{\le}{M}$ iku transitif. Mula, model-model saka
\begin{align*}
\{\quad &\lforall[x][x \le x], \\
   & \lforall[x][\lforall[y][((x \le y \land y \le
    x) \lif x = y)]], \\
   &\lforall[x][\lforall[y][\lforall[z][((x \le y
      \land y \le z) \lif x \le z)]]] \quad \}
\end{align*}
iku persis struktur-struktur kang~$\Assign{\le}{M}$ refleksif,
antisimetris, lan transitif, yaiku urutan parsial. Mula,
kita bisa nggunakake ukara-ukara iku minangka aksioma
kanggo \emph{teori logika tataran kapisan babagan urutan parsial}.
\end{ex}

\end{document}
